\section{Giới thiệu}
\label{sec:introduction}

\subsection{Tổng quan về đề tài}
\label{ssec:overview}

Đề tài này tập trung vào việc thiết kế và triển khai một ứng dụng web quản lý công việc (Task Manager) sử dụng kiến trúc serverless hoàn toàn trên nền tảng Amazon Web Services (AWS). Mục tiêu chính của đề tài là xây dựng một hệ thống có khả năng tự động mở rộng (auto-scaling), đảm bảo tính sẵn sàng cao (High Availability), bảo mật theo nguyên tắc phòng thủ nhiều lớp (defense-in-depth), và quan trọng nhất là hoạt động hoàn toàn trong giới hạn miễn phí (Free Tier) của AWS.

Kiến trúc serverless đã trở thành xu hướng phát triển ứng dụng hiện đại nhờ khả năng loại bỏ gánh nặng quản lý hạ tầng, tự động mở rộng theo nhu cầu thực tế, và mô hình thanh toán theo mức sử dụng (pay-per-use). Trong đề tài này, nhóm không sử dụng bất kỳ máy ảo (EC2) nào, thay vào đó tận dụng các dịch vụ được quản lý hoàn toàn (fully managed services) như AWS Lambda, API Gateway, DynamoDB, CloudFront và S3.

\subsection{Mục tiêu của đề tài}
\label{ssec:objectives}

Đề tài hướng đến việc đạt được các mục tiêu sau:

\begin{enumerate}
    \item \textbf{Kiến trúc 100\% serverless:} Toàn bộ hệ thống không sử dụng máy ảo, chỉ dựa vào các dịch vụ serverless của AWS.
    
    \item \textbf{Tự động mở rộng:} Hệ thống tự động điều chỉnh tài nguyên dựa trên lưu lượng truy cập thực tế mà không cần cấu hình thủ công.
    
    \item \textbf{Tính sẵn sàng cao:} Thiết kế hệ thống trải rộng trên nhiều Availability Zone (AZ) để đảm bảo hoạt động liên tục ngay cả khi một AZ gặp sự cố.
    
    \item \textbf{Cách ly mạng:} Lambda kết nối với DynamoDB hoàn toàn qua mạng riêng của AWS (VPC Endpoint), không đi qua Internet công cộng, đảm bảo bảo mật và không phát sinh chi phí NAT Gateway.
    
    \item \textbf{Bảo mật frontend:} S3 bucket lưu trữ frontend phải ở chế độ private, chỉ có thể truy cập qua CloudFront với Origin Access Control (OAC).
    
    \item \textbf{Chi phí bằng 0:} Toàn bộ hệ thống hoạt động trong giới hạn Free Tier của AWS, không phát sinh chi phí.
    
    \item \textbf{Giám sát và kiểm soát chi phí:} Triển khai CloudWatch Dashboard, Alarms và AWS Budgets để theo dõi hoạt động và kiểm soát chi phí.
\end{enumerate}

\subsection{Phạm vi thực hiện}
\label{ssec:scope}

Đề tài bao gồm các phần chính sau:

\begin{itemize}
    \item \textbf{Frontend:} Ứng dụng web đơn giản sử dụng HTML, CSS, JavaScript thuần, cung cấp giao diện quản lý công việc với đầy đủ chức năng CRUD (Create, Read, Update, Delete) và khả năng lọc theo ngày đến hạn và mức độ ưu tiên.
    
    \item \textbf{Backend API:} Xây dựng RESTful API sử dụng AWS Lambda (Python 3.12) và API Gateway, cung cấp 4 endpoints chính: GET /tasks, POST /tasks, PUT /tasks/:id, DELETE /tasks/:id.
    
    \item \textbf{Database:} Sử dụng Amazon DynamoDB (NoSQL) với Global Secondary Index (GSI) để truy vấn hiệu quả theo userId.
    
    \item \textbf{Networking:} Thiết lập Custom VPC với 2 Private Subnets trải rộng trên 2 AZ, cấu hình VPC Gateway Endpoint cho DynamoDB để đảm bảo kết nối private.
    
    \item \textbf{Security:} Áp dụng nguyên tắc Least Privilege cho IAM Roles, cấu hình S3 private với CloudFront OAC, và Security Group hạn chế chỉ cho phép kết nối đến DynamoDB.
    
    \item \textbf{Monitoring:} Triển khai CloudWatch Dashboard với các metrics quan trọng, cấu hình Alarms gửi thông báo qua SNS, và thiết lập AWS Budgets với ngưỡng \$0.01.
\end{itemize}

\subsection{Cấu trúc báo cáo}
\label{ssec:structure}

Báo cáo được tổ chức thành các phần chính như sau:

\begin{itemize}
    \item \textbf{Phần \ref{sec:architecture}:} Trình bày chi tiết về thiết kế kiến trúc hệ thống, giải thích luồng xử lý request và cơ chế auto-scaling.
    
    \item \textbf{Phần \ref{sec:implementation}:} Mô tả quá trình triển khai từng thành phần của hệ thống bao gồm Frontend, Backend, Database, Networking và IAM.
    
    \item \textbf{Phần \ref{sec:security}:} Cung cấp bằng chứng về các biện pháp bảo mật đã được triển khai thông qua screenshots và giải thích chi tiết.
    
    \item \textbf{Phần \ref{sec:monitoring}:} Trình bày hệ thống giám sát, cảnh báo và kiểm soát chi phí.
    
    \item \textbf{Phụ lục:} Hướng dẫn chi tiết về cách triển khai và source code.
\end{itemize}
